\lstdefinelanguage{clingo}{
  morecomment=[l]{\%},
  morekeywords={not},
  sensitive=true
}
\lstdefinestyle{clingostyle}{
  language=clingo,
  basicstyle=\ttfamily\footnotesize,
  commentstyle=\color{gray}\itshape,
  keywordstyle=\color{blue!60!black},
  columns=fullflexible,
  keepspaces=true,
  breaklines=true,
  frame=single,
  framerule=0.3pt,
  rulecolor=\color{black!30},
  xleftmargin=2pt,
  xrightmargin=2pt,
  aboveskip=0.8em,
  belowskip=0.6em
}
\lstset{style=clingostyle}

%%%%%%%%%%%%%%%%%%%%%%%%%%%%%%%%%%%%%%%%%%%%%%%%%%5
% TAKE THIS SECTION AS AN EXAMPLE TO ASP ENCODING, THE REAL ASP PROGRAM IS STILL BEING TESTED.
%%%%%%%%%%%%%%%%%%%%%%%%%%%%%%%%%%%%%%%%%%%%%%%%%%
\subsection{ASP Encoding}
\label{sec:encoding}
% =====================================================================
 
\subsubsection{Design}
 
The formulation is encoded in the answer set programming language of
\texttt{clingo}. We use ASP, rather than a generic CSP library, because
sums, counting, optimization with priorities and conditional (``exists a
clause member'') reasoning are built into the language.
The program has the usual \emph{generate and test} shape:
a choice rule generates the free cells $x_{c,s}$, and the constraints reject
plans.
 
Three design decisions are worth stating.
\begin{enumerate}
  \item \textbf{Constraints are violation atoms.}
        Instead of writing each constraint as an integrity constraint, we
        derive an atom \texttt{viol(kind,C,S)} when it is broken, and add a
        single integrity constraint forbidding all violations. If the
        external atom \texttt{check\_mode} is true, the integrity constraint
        is switched off and the program reports the violations. This gives
        the check mode without a second program.
  \item \textbf{One program, three modes.} Generate mode needs no draft.
        Repair mode adds the draft as facts \texttt{prop(C,S)}, and the
        distance to it is the highest-priority objective. Check mode adds
        \texttt{check\_mode} and forces distance $0$.
  \item \textbf{Integers only.} \texttt{clingo} has no floating-point
        numbers, so grades are stored in points $\times 100$
        ($3.00\to300$) and the GPA test is written by cross-multiplication:
        $\sum \mathrm{cr}\cdot\gamma \ge \theta \sum \mathrm{cr}$.
\end{enumerate}
 
The encoding is split into three files: \texttt{core.lp} (the formulation),
\texttt{data.lp} (the curriculum) and \texttt{student.lp} (history, expected
grades, budget). Table~\ref{tab:mapping} shows where each part of the
formulation lives.
 
\begin{table}[ht]
\centering
\small
\begin{tabular}{@{}p{4.2cm}>{\raggedright\arraybackslash}p{8.8cm}@{}}
\toprule
\textbf{Formulation} & \textbf{Encoding} \\
\midrule
Variable $x_{c,s}$ & choice rule \texttt{take(C,S)} for $s>s_0$; history in \texttt{hist(C,S,G)} \\
Pass, $P^{<}$, $P^{\le}$ & \texttt{pass}, \texttt{passed\_before}, \texttt{passed\_upto} \\
Load $L_s$, earned $E_s$, level & \texttt{load}, \texttt{earned}, \texttt{earned\_before}, \texttt{lam} \\
C1--C7 & atoms \texttt{viol(offering,\dots)}, \texttt{viol(retake,\dots)}, \texttt{viol(prereq,\dots)}, \texttt{viol(coreq,\dots)}, \texttt{viol(maxload,\dots)}, \texttt{viol(minload,\dots)}, \texttt{viol(fee,\dots)}, \texttt{viol(level,\dots)} \\
G1--G3, $T$ & \texttt{gradok}, \texttt{first\_grad}; \texttt{viol(no\_grad,\dots)} \\
G4 (GPA) & \texttt{best}, \texttt{gpa\_pts}, \texttt{gpa\_cr}, \texttt{short} \\
Objectives & \texttt{\#minimize} statements with priorities \\
\bottomrule
\end{tabular}
\caption{Mapping from the formulation to the encoding.}
\label{tab:mapping}
\end{table}
 
\subsubsection{Core Program}
 
\paragraph{Parameters and variables}
Parameters are constants that can be overridden from the command line.
The fee function is a table from load to money.
\begin{lstlisting}
% ===== Parameters (overridable with -c name=value) =====
#const now       = 1.     % last completed semester
#const tmax      = 5.     % planning horizon
#const lmin      = 6.     % minimum load (credits)
#const lmax      = 9.     % maximum load (credits)
#const lcap      = 30.    % size of the fee table
#const passmark  = 100.   % passing grade, in points x 100
#const credreq   = 24.    % credits required to graduate
#const gpamin    = 250.   % required cumulative GPA (2.50 x 100)
#const feeunit   = 1000.  % fee per credit (example fee function)
#external check_mode.     % true => only report violations
#defined ext_req/1. #defined ext_ok/1. #defined improve/1.
#defined fixed/2. #defined banned/2. #defined prop/2.
 
sem(1..tmax).
fee(L, feeunit*L) :- L = 0..lcap.      % fee function: load -> money
\end{lstlisting}
\begin{lstlisting}
% ===== Variable: planned cells take(C,S) in {0,1} =====
{ take(C,S) : course(C), sem(S), S > now }.
\end{lstlisting}
 
\paragraph{Attempts, passes and credits}
A planned attempt passes by assumption A1 (\texttt{pass(C,S) :- take(C,S)}),
and only the first pass of a course is earned.
\begin{lstlisting}
% ===== Attempts, grades, passes =====
att(C,S)      :- hist(C,S,_).
att(C,S)      :- take(C,S).
grade(C,S,G)  :- hist(C,S,G).
grade(C,S,G)  :- take(C,S), exp(C,G).
pass(C,S)     :- hist(C,S,G), G >= passmark.
pass(C,S)     :- take(C,S).                       % optimistic pass
 
passed_before(C,S) :- pass(C,S0), sem(S), S0 < S.
passed_upto(C,S)   :- pass(C,S0), sem(S), S0 <= S.
hist_passed(C)     :- hist(C,S,G), G >= passmark.
plan_passed_before(C,S) :- take(C,S0), S0 > now, sem(S), S0 < S.
attempt_before(C,S)     :- att(C,S0), sem(S), S0 < S.
\end{lstlisting}
\begin{lstlisting}
% ===== Load, earned credits, level =====
load(S,L)  :- sem(S), L = #sum{ Cr,C : att(C,S), credits(C,Cr) }.
earned(C,S) :- pass(C,S), not passed_before(C,S).  % first pass only
earned_before(S,E) :- sem(S),
    E = #sum{ Cr,C : earned(C,S0), S0 < S, credits(C,Cr) }.
earned_upto(S,E)   :- sem(S),
    E = #sum{ Cr,C : earned(C,S0), S0 <= S, credits(C,Cr) }.
lam(S,L) :- earned_before(S,E), lvl_range(L,Lo,Hi), Lo <= E, E < Hi.
\end{lstlisting}
 
\paragraph{Graduation and GPA}
\texttt{first\_grad(T)} is the graduation semester. The retake policy
``highest grade counts'' is written with \texttt{better} and \texttt{best};
the policy ``latest grade counts'' only changes these two rules.
\begin{lstlisting}
% ===== Graduation =====
miss_comp(S) :- sem(S), comp(C), not passed_upto(C,S).
miss_grp(S)  :- sem(S), reqgrp(G,R),
    #sum{ Cr,C : earned(C,S0), S0 <= S, ingroup(C,G), credits(C,Cr) } < R.
ext_fail     :- ext_req(X), not ext_ok(X).
gradok(S)    :- earned_upto(S,E), E >= credreq,
                not miss_comp(S), not miss_grp(S), not ext_fail.
first_grad(T) :- gradok(T), not gradok(S) : sem(S), S < T.
\end{lstlisting}
\begin{lstlisting}
% ===== GPA (retake policy: highest grade counts) =====
better(C,G) :- grade(C,_,G), grade(C,_,G2), G2 > G.
best(C,G)  :- grade(C,_,G), not better(C,G).
gpa_pts(P) :- P = #sum{ G*Cr,C : best(C,G), credits(C,Cr) }.
gpa_cr(Q)  :- Q = #sum{ Cr,C : best(C,_), credits(C,Cr) }.
short(D)   :- gpa_pts(P), gpa_cr(Q), D = gpamin*Q - P, D > 0.
\end{lstlisting}
 
\paragraph{Constraints}
Prerequisites and corequisites are given in CNF as facts
\texttt{pre(C,K,A)}: course \texttt{C} has clause \texttt{K} containing
alternative \texttt{A}. A clause is satisfied when some alternative is
passed in time (\texttt{presat}, \texttt{cosat}).
\begin{lstlisting}
% ===== Constraints, written as violation atoms =====
presat(C,K,S) :- pre(C,K,A), passed_before(A,S).
cosat(C,K,S)  :- co(C,K,A),  passed_upto(A,S).
 
viol(offering,  C,S) :- take(C,S), not offered(C,S).
viol(retake,    C,S) :- take(C,S), plan_passed_before(C,S).
viol(retake,    C,S) :- take(C,S), hist_passed(C), not improve(C).
viol(prereq,    C,S) :- take(C,S), pre(C,K,_), not presat(C,K,S).
viol(coreq,     C,S) :- take(C,S), co(C,K,_),  not cosat(C,K,S).
viol(level,     C,S) :- take(C,S), level(C,V), lam(S,L), V > L.
viol(maxload,   0,S) :- load(S,L), S > now, L > lmax.
viol(minload,   0,S) :- load(S,L), S > now, first_grad(T), S < T, L < lmin.
viol(fee,       0,S) :- load(S,L), S > now, fee(L,F), budget(S,B), F > B.
viol(after_grad,C,S) :- take(C,S), first_grad(T), S > T.
viol(no_grad,   0,0) :- #count{ T : first_grad(T) } = 0.
 
:- viol(_,_,_), not check_mode.
\end{lstlisting}
\begin{lstlisting}
% ===== Student decisions =====
:- fixed(C,S), not take(C,S).
:- banned(C,S), take(C,S).
\end{lstlisting}
 
\paragraph{Repair distance and objectives}
In \texttt{clingo}, a larger priority number is more important, so the
distance to the draft has the highest priority (6) and the load balance the
lowest (2).
\begin{lstlisting}
% ===== Distance to a proposed plan (repair mode) =====
has_prop :- prop(_,_).
changed(C,S) :- prop(C,S), not take(C,S).
changed(C,S) :- take(C,S), not prop(C,S), has_prop.
:- changed(_,_), check_mode.
\end{lstlisting}
\begin{lstlisting}
% ===== Objectives (higher priority number = more important) =====
maxload(M) :- M = #max{ L : load(S,L), S > now }.
#minimize{ 1@6,C,S : changed(C,S) }.
#minimize{ T@5     : first_grad(T) }.
#minimize{ D@4     : short(D) }.
#minimize{ 1@3,C,S : take(C,S), attempt_before(C,S) }.
#minimize{ M@2     : maxload(M) }.
\end{lstlisting}
 
\subsubsection{Interaction}
 
The three modes are selected only by the input files:
\begin{lstlisting}[language={},basicstyle=\ttfamily\footnotesize]
clingo core.lp data.lp student.lp                         # generate
clingo core.lp data.lp student.lp draft.lp check.lp       # check
clingo core.lp data.lp student.lp draft.lp                # repair
\end{lstlisting}
where \texttt{draft.lp} lists the student's draft as \texttt{prop(C,S)} facts
and \texttt{check.lp} contains the single fact \texttt{check\_mode.}
Pins and bans are facts \texttt{fixed(C,S)} and \texttt{banned(C,S)}.
For a dialogue with the student, the same program can be driven from the
\texttt{clingo} Python API, which can toggle such facts as external atoms and
re-solve without grounding everything again.
 
% ===================================================================== The below section is NOT VERIFIED YET
% \subsection{Worked Example}
% \label{sec:example}
% % =====================================================================
 
% We run the encoding on a toy curriculum small enough to verify by hand.
% All results below were produced with \texttt{clingo} 5.8 and agree with the
% manual analysis.
 
% \subsubsection{Setting}
 
% \paragraph{Curriculum}
% Table~\ref{tab:toycourses} lists nine courses of $3$ credits each.
% The parameters are: $L_{\min}=6$, $L_{\max}=9$, $\mathrm{Cr}_{\min}=24$,
% $\theta=2.50$, $T_{\max}=5$, and $\mathrm{fee}(L)=1000\cdot L$.
% The level function is $\lambda(E)=1$ for $E<6$, $2$ for $6\le E<12$ and $3$
% for $E\ge 12$. At least $6$ credits must come from the elective group
% \{EL1, EL2, EL3\}. External requirements are assumed satisfied.
 
% \begin{table}[ht]
% \centering
% \small
% \begin{tabular}{@{}lccp{3.2cm}>{\raggedright\arraybackslash}p{4.4cm}@{}}
% \toprule
% \textbf{Course} & \textbf{Level} & \textbf{Compulsory} & \textbf{Offered} & \textbf{Relations} \\
% \midrule
% MA101 & 1 & yes & every semester & --- \\
% MA102 & 1 & no  & odd semesters  & pre: MA101 \\
% CS101 & 1 & yes & every semester & --- \\
% CS102 & 2 & yes & every semester & pre: CS101 $\wedge$ (MA101 $\vee$ MA102) \\
% CS201 & 2 & yes & every semester & pre: CS101; co: CS102 \\
% CS301 & 3 & yes & even semesters & pre: CS201 \\
% EL1, EL2 & 1 & no & every semester & --- \\
% EL3 & 2 & no & every semester & --- \\
% \bottomrule
% \end{tabular}
% \caption{Toy curriculum (all courses have 3 credits).}
% \label{tab:toycourses}
% \end{table}
 
% \paragraph{Student}
% The student has finished semester $1$ ($s_0=1$): MA101 passed with a B
% ($3.00$) and CS101 \emph{failed}, so CS101 must be retaken. The load of
% semester $1$ is $6$ credits but only $3$ credits are earned. Every planned
% course is expected to get a B.
 
% \subsubsection{Scenario 1: generating a plan}
 
% With a budget of $9000$ per semester (so the fee allows the full $9$
% credits), the solver returns the plan in Table~\ref{tab:plan1}, found as
% optimal with $T=4$. Several plans share the same optimal values; the
% solver returns one of them.
 
% \begin{table}[ht]
% \centering
% \small
% \begin{tabular}{@{}lccccc@{}}
% \toprule
% \textbf{Course} & \textbf{S1 (history)} & \textbf{S2} & \textbf{S3} & \textbf{S4} & \textbf{S5} \\
% \midrule
% MA101 & passed &        &        &        &  \\
% CS101 & failed & $\bullet$ &        &        &  \\
% EL1   &        & $\bullet$ &        &        &  \\
% MA102 &        &        & $\bullet$ &        &  \\
% CS102 &        &        & $\bullet$ &        &  \\
% CS201 &        &        & $\bullet$ &        &  \\
% CS301 &        &        &        & $\bullet$ &  \\
% EL2   &        &        &        & $\bullet$ &  \\
% \midrule
% \textbf{Load} & 6 & 6 & 9 & 6 & 0 \\
% \bottomrule
% \end{tabular}
% \caption{Optimal plan with budget $9000$ ($\bullet$ = planned attempt).}
% \label{tab:plan1}
% \end{table}
 
% The plan follows the constraints as expected. CS101 is retaken first
% because CS102 and CS201 depend on it. CS102 and CS201 share semester $3$,
% which is allowed because CS102 is a corequisite of CS201. CS301 is placed in
% semester $4$, the first even semester in which its prerequisite is passed.
% The earned credits are $3+21=24=\mathrm{Cr}_{\min}$, the elective group
% gets $6$ credits, one planned retake is used and the projected GPA is
% $3.00\ge 2.50$. The minimum load is waived in the graduation semester,
% and semester $5$ is empty.
 
% \subsubsection{Scenario 2: a tighter budget}
 
% If the budget drops to $6000$ per semester, the fee function limits the load
% to $6$ credits ($\mathrm{fee}(9)=9000>6000$). The chain CS101 $\to$ CS102/CS201
% $\to$ CS301 is unchanged, but the remaining $21$ credits no longer fit in
% three semesters, so the optimal plan graduates in semester $5$:
% S2 = \{CS101, EL1\}, S3 = \{CS102, CS201\}, S4 = \{CS301, EL2\} and
% S5 = \{MA102\}. MA102 is placed in semester $5$ because it is offered only
% in odd semesters, and its load of $3$ is below $L_{\min}$, which is allowed
% in the graduation semester.
 
% \subsubsection{Scenario 3: checking and repairing a draft}
 
% Suppose the student proposes, with the budget of $9000$:
% \[
%   \text{S2} = \{\text{CS101},\ \text{CS102},\ \text{EL1}\}, \qquad
%   \text{S3} = \{\text{CS201},\ \text{CS301},\ \text{EL2}\}.
% \]
 
% \paragraph{Check}
% The program reports the violations in Table~\ref{tab:viol}.
 
% \begin{table}[ht]
% \centering
% \small
% \begin{tabular}{@{}lp{9.5cm}@{}}
% \toprule
% \textbf{Violation} & \textbf{Reason} \\
% \midrule
% \texttt{prereq(cs102,2)} & CS101 is not yet passed before semester 2 (it was failed and is retaken in the same semester) \\
% \texttt{level(cs102,2)} & CS102 has level 2, but the student has earned only $3$ credits before semester 2 (level 1) \\
% \texttt{offering(cs301,3)} & CS301 is not offered in odd semesters \\
% \texttt{prereq(cs301,3)} & CS201 is planned in the same semester, not an earlier one \\
% \texttt{no\_grad} & the draft earns only $3+18=21<24$ credits and does not graduate within the horizon \\
% \bottomrule
% \end{tabular}
% \caption{Violations reported by the check mode.}
% \label{tab:viol}
% \end{table}
 
% \paragraph{Repair}
% The repair mode returns the feasible plan with the fewest changed cells:
% S2 = \{CS101, EL1\}, S3 = \{CS102, CS201, EL2\}, S4 = \{CS301, EL3\}, with
% $T=4$. Compared with the draft it changes $5$ cells:
% \begin{itemize}
%   \item move CS102 from semester 2 to semester 3 (2 cells);
%   \item move CS301 from semester 3 to semester 4 (2 cells);
%   \item add EL3 in semester 4 (1 cell), to reach the required credits.
% \end{itemize}
% These changes are the explanation given to the student. EL3 is chosen
% rather than MA102 because MA102 is offered only in odd semesters and
% semester $3$ is already full. If the budget of $6000$ of Scenario~2 applies to the
% same draft, the minimum repair needs $7$ changed cells and graduation moves
% to semester~$5$.
 
% \subsubsection{Summary}
 
% The example shows each part of the formulation at work: the retake counted
% in the load but not in the earned credits, CNF prerequisites, a corequisite
% in the same semester, the level and offering restrictions, the fee limiting
% the load, and the check and repair modes producing an explanation instead
% of a bare ``no plan''.